\documentclass[12pt,a4paper]{article}

% ── Packages ──────────────────────────────────────────────────────────────────
\usepackage[T1]{fontenc}
\usepackage[utf8]{inputenc}
\usepackage{graphicx}
\usepackage{geometry}
\usepackage{hyperref}
\usepackage{booktabs}
\usepackage{tabularx}
\usepackage{longtable}
\usepackage{array}
\usepackage{enumitem}
\usepackage{xcolor}
\usepackage{fancyhdr}
\usepackage{titlesec}
\usepackage{parskip}
\usepackage{microtype}
\usepackage{listings}
\usepackage{pifont}   % for \ding checkmarks
\usepackage{mdframed}
\usepackage{setspace}
\usepackage{tocloft}
\usepackage{float}
\usepackage{caption}

% ── Page geometry ─────────────────────────────────────────────────────────────
\geometry{
  top=2.5cm, bottom=2.5cm,
  left=2.5cm, right=2.5cm
}

% ── Colours ───────────────────────────────────────────────────────────────────
\definecolor{primary}{RGB}{37,99,235}      % blue-600
\definecolor{accent}{RGB}{16,185,129}      % emerald-500
\definecolor{lightgray}{RGB}{245,245,245}
\definecolor{darkgray}{RGB}{55,65,81}
\definecolor{codebg}{RGB}{30,30,30}
\definecolor{codegreen}{RGB}{166,226,46}

% ── Hyperref setup ────────────────────────────────────────────────────────────
\hypersetup{
  colorlinks=true,
  linkcolor=primary,
  urlcolor=primary,
  citecolor=primary,
  pdftitle={QuizPod -- Progress-1 Documentation},
  pdfauthor={Dev Aggarwal et al.}
}

% ── Section title formatting ──────────────────────────────────────────────────
\titleformat{\section}
  {\Large\bfseries\color{primary}}
  {\thesection.}{0.5em}{}[\vspace{-0.3em}\textcolor{primary}{\rule{\linewidth}{1.5pt}}]

\titleformat{\subsection}
  {\large\bfseries\color{darkgray}}
  {\thesubsection}{0.5em}{}

\titleformat{\subsubsection}
  {\normalsize\bfseries\color{darkgray}}
  {\thesubsubsection}{0.5em}{}

% ── Header / footer ───────────────────────────────────────────────────────────
\pagestyle{fancy}
\fancyhf{}
\fancyhead[L]{\small\textcolor{darkgray}{QuizPod --- Web Technology Progress-1}}
\fancyhead[R]{\small\textcolor{darkgray}{March 2026}}
\fancyfoot[C]{\small\thepage}
\renewcommand{\headrulewidth}{0.4pt}
\renewcommand{\footrulewidth}{0pt}

% ── Code listing style ────────────────────────────────────────────────────────
\lstset{
  basicstyle=\ttfamily\small,
  backgroundcolor=\color{lightgray},
  frame=single,
  rulecolor=\color{gray!40},
  breaklines=true,
  columns=fullflexible,
  keepspaces=true,
  showstringspaces=false,
  tabsize=2,
}

% ── Convenience commands ──────────────────────────────────────────────────────
\newcommand{\yes}{\textcolor{accent}{\ding{51}}}   % tick
\newcommand{\no}{\textcolor{red}{\ding{55}}}        % cross

% ── Highlighted box ───────────────────────────────────────────────────────────
\newmdenv[
  backgroundcolor=blue!5,
  linecolor=primary,
  linewidth=1.5pt,
  roundcorner=4pt,
  skipabove=8pt,skipbelow=8pt,
  innerleftmargin=10pt,innerrightmargin=10pt,
  innertopmargin=8pt,innerbottommargin=8pt
]{infobox}

\graphicspath{{./}}

% ─────────────────────────────────────────────────────────────────────────────
%  TITLE PAGE
% ─────────────────────────────────────────────────────────────────────────────
\title{
  \vspace{-1cm}
  {\LARGE\bfseries\color{primary} QuizPod}\\[0.4em]
  {\large An Intelligent Quiz \& Learning Platform}\\[0.8em]
  {\normalsize\color{darkgray} Web Technology --- Progress-1 Documentation}
}
\author{
  \begin{tabular}{lr}
    \textbf{Name} & \textbf{Roll No.} \\[0.4em]
    Ayush Chauhan  & 2023UIT3134 \\
    Dev Aggarwal   & 2023UIT3136 \\
    Bhavya Agrawal & 2023UIT3138 \\
    Rohan Singh    & 2023UIT3139 \\
    Payal          & 2023UIT3142 \\
    Mudit          & 2023UIT3149 \\
  \end{tabular}\\[0.6em]
  {\small Team Size: 6 Members}
}
\date{March 2026}

% ─────────────────────────────────────────────────────────────────────────────
\begin{document}
% ─────────────────────────────────────────────────────────────────────────────

\maketitle
\thispagestyle{empty}

\begin{infobox}
  \textbf{Subject:} Web Technology \hfill
  \textbf{Evaluation:} Progress-1\\
  \textbf{GitHub:} \href{https://github.com/DevAggarwal03/WebTech_Project}{\texttt{DevAggarwal03/WebTech\_Project}}
\end{infobox}

\vspace{0.5cm}
\tableofcontents
\newpage

% =============================================================================
\section{Topic of the Project}
% =============================================================================

\subsection{Problem Statement and Theme}

Existing online quiz and learning platforms fail to genuinely optimise for
\emph{learning outcomes}. They prioritise engagement metrics (leaderboards,
speed-based scoring) or administrative convenience, while students receive
shallow feedback, no adaptive learning, and teachers get minimal actionable
analytics.

\textbf{Core Problem Areas:}

\begin{itemize}[leftmargin=1.5em]
  \item \textbf{Shallow feedback} --- Platforms show scores but never explain
        \emph{why} an answer was wrong or \emph{what to study next}.
  \item \textbf{No adaptive difficulty} --- Every student receives the same
        quiz regardless of skill level.
  \item \textbf{Academic integrity gaps} --- Tab-switching, answer-sharing, and
        unauthorised quiz access are poorly addressed.
  \item \textbf{Unauthorised access} --- Students not enrolled in a specific
        class can freely join if they obtain the access code, compromising quiz
        integrity.
  \item \textbf{Paywall exclusion} --- Analytics, student tracking, and rich
        question types are locked behind expensive paid tiers.
  \item \textbf{No collaborative learning} --- Quizzes are always competitive;
        team-based learning modes are absent.
\end{itemize}

\textbf{Our Solution --- QuizPod:}\\
A web-based quiz platform that places learning outcomes at the centre. It
provides rich feedback, teacher-controlled classrooms with unauthorised-student
detection, detailed performance analytics, and lightweight engagement elements
like XP and badges --- all for free.

\begin{infobox}
  \textit{``The quiz platform that actually helps students learn, not just
  score --- and actually helps teachers teach, not just grade.''}
\end{infobox}

% -----------------------------------------------------------------------------
\subsection{Research on Existing Platforms}

We conducted detailed research on six major quiz/assessment platforms to
identify gaps and opportunities.

\begin{table}[H]
  \centering
  \caption{Analysis of Existing Quiz Platforms}
  \small
  \begin{tabularx}{\linewidth}{@{} l l l X @{}}
    \toprule
    \textbf{Platform} & \textbf{Primary Strength} & \textbf{Core Audience} & \textbf{Key Shortcoming} \\
    \midrule
    Kahoot       & Gamified live quizzes    & K-12, higher ed   & Speed $>$ accuracy; free tier capped at 10 players \\
    Quizlet      & Flashcards \& study modes & Students         & Key modes behind paywall; ad-heavy free tier \\
    Google Forms & Simple form-based quizzes & General          & No native timer; easy misconfiguration exposes answers \\
    Mentimeter   & Real-time polling        & Higher ed         & Only 2 free interactive slides; weak analytics \\
    Socrative    & Formative live assessments & K-12 teachers   & Capped at 50 students; MCQ-only \\
    Typeform     & Conversational form UX   & Business/HR       & Expensive per-response pricing; no auto-close \\
    \bottomrule
  \end{tabularx}
\end{table}

\subsubsection{Universal Gaps Across All Platforms}

\begin{enumerate}[leftmargin=1.5em]
  \item No Adaptive Difficulty --- questions never adjust based on student performance.
  \item Shallow Feedback --- most platforms show a score with no explanation.
  \item No Spaced Repetition --- no intelligent review scheduling.
  \item Weak Academic Integrity --- no lightweight anti-cheat mechanisms.
  \item Paywalls block core features --- analytics, tracking, and rich types are premium-only.
  \item No Collaborative Quizzing --- students always compete, never cooperate.
  \item Poor Teacher Analytics --- class averages but not per-concept diagnosis.
  \item Rich Content Limitations --- code, \LaTeX{}, diagrams not natively supported.
  \item Fragmented Workflow --- quiz creation, sharing, and grading are spread across tools.
  \item No Offline Mode --- all platforms require stable internet.
\end{enumerate}

\subsubsection{Research on Unauthorised Access Prevention}

\begin{table}[H]
  \centering
  \caption{Access Control Mechanisms Across Platforms}
  \small
  \begin{tabularx}{\linewidth}{@{} l l X @{}}
    \toprule
    \textbf{Platform} & \textbf{Access Mechanism} & \textbf{Gaps} \\
    \midrule
    Google Classroom & 6--8 char class code; domain restriction & Code can be shared/leaked; no approval step \\
    Kahoot           & Temporary game PIN per session           & No persistent ``class'' concept \\
    Quizizz          & Join code via joinmyquiz.com            & Code-based, no approval \\
    Moodle           & Enrollment key, manual enrollment       & Powerful but complex for teachers \\
    Canvas LMS       & SIS integration, self-enroll codes      & Institution-managed; limited teacher control \\
    Edmodo           & Group code + teacher approval queue     & \yes Best approach --- but shut down in 2022 \\
    \bottomrule
  \end{tabularx}
\end{table}

\textbf{Key Insight:} Most platforms rely on \emph{security by obscurity} --- a
shared code is the only gate. Codes get leaked via social media, there is no
identity verification, and no teacher-approval step in popular quiz tools.
Edmodo's approval-queue approach was the best but is no longer available.

\textbf{Our three-part approach:}
\begin{enumerate}[leftmargin=1.5em]
  \item \textbf{Teacher approval queue} --- students request to join; teachers approve/deny.
  \item \textbf{Class roster matching} --- students enter their roll number, matched against a pre-uploaded list.
  \item \textbf{Auto-expiring join codes} --- codes valid only for a limited time window.
\end{enumerate}

\begin{infobox}
  This combination does not exist in any current lightweight quiz platform.
\end{infobox}

% =============================================================================
\section{Structure of the Project with Technology Stack}
% =============================================================================

\subsection{Project Structure --- Pages and Sitemap}

\begin{lstlisting}
QuizPod Platform
|
+-- Landing Page (/)
|   +-- Hero section with tagline & CTA buttons (Sign Up / Join Quiz)
|   +-- Feature cards (Practice Mode, Live Competitions, Track Progress)
|   `-- Footer with social links
|
+-- Authentication (/auth)
|   +-- Role selection toggle (Student / Teacher)
|   +-- Sign-up form (Name, Email, Password)
|   `-- Animated binary-rain background
|
+-- Student Dashboard (/student)
|   +-- Navbar (Home, XP counter, Join Quiz, Logout)
|   +-- Hero banner with stats (Rank, XP)
|   +-- Badge showcase
|   `-- Recent Quizzes grid (title, score, date, category)
|
+-- Teacher Dashboard (/teacher)
|   +-- Navbar (Home, Community, Create Quiz, Logout)
|   +-- Hero banner with stats (Total Quizzes, Upcoming Quizzes)
|   +-- Subject badges (Science, Math, English, Geography)
|   `-- Recent Quizzes grid (title, subject, submissions, date)
|
+-- Quiz Arena (/quiz/:id)
|   +-- Navbar (brand, section name, theme toggle, countdown timer)
|   +-- Question panel (text, MCQ radio options)
|   +-- Navigation (Previous, Mark for Review, Next/Submit)
|   +-- Status panel (colour-coded question grid)
|   `-- Built-in calculator tool
|
+-- Create Quiz (/teacher/create-quiz)
|   `-- Quiz builder for teachers
|
`-- Create Classroom (/teacher/create-classroom)
    `-- Classroom creation, join codes, roster, enrolment management
\end{lstlisting}

% ── Frontend screenshot placeholders ─────────────────────────────────────────
\begin{figure}[H]
  \centering
  \includegraphics[width=0.47\linewidth]{assets/landing.png}
  \hfill
  \includegraphics[width=0.47\linewidth]{assets/auth.png}
  \caption{Frontend Screenshots --- Landing Page (left) and Authentication Page (right)}
\end{figure}

\begin{figure}[H]
  \centering
  \includegraphics[width=0.47\linewidth]{assets/studentDashboard.jpeg}
  \hfill
  \includegraphics[width=0.47\linewidth]{assets/TeacherDashboard.jpeg}
  \caption{Frontend Screenshots --- Student Dashboard (left) and Teacher Dashboard (right)}
\end{figure}

\begin{figure}[H]
  \centering
  \includegraphics[width=0.47\linewidth]{assets/QuizArena.png}
  \caption{Frontend Screenshot --- Quiz Arena}
\end{figure}

\begin{table}[H]
  \centering
  \caption{Design Principles Applied in QuizPod}
  \small
  \begin{tabularx}{\linewidth}{@{} l X @{}}
    \toprule
    \textbf{Principle} & \textbf{Application in QuizPod} \\
    \midrule
    Feedback-first, score-second    & Prioritise explaining wrong answers over raw scores \\
    No public shaming               & Private leaderboards; personal stats on dashboard \\
    Teacher time is precious        & Streamlined dashboard, at-a-glance stats \\
    Mobile-first approach           & Responsive CSS for all screen sizes \\
    Component-based architecture    & Each page is a self-contained React component with its own stylesheet \\
    \bottomrule
  \end{tabularx}
\end{table}

\textbf{Design Topology:}
\begin{itemize}[leftmargin=1.5em]
  \item \textbf{Component-based architecture} using React --- modular, reusable UI components for consistency.
  \item \textbf{Colour-coded status system} in Quiz Arena: \textcolor{accent}{green} = answered, \textcolor{orange}{yellow} = marked for review, \textcolor{gray}{grey} = not visited.
  \item \textbf{Role-based interface separation} --- Student and Teacher dashboards are entirely different views with distinct layouts, stats, and actions.
  \item \textbf{Consistent visual language} --- shared design tokens (colours, spacing, typography) across all pages.
\end{itemize}

% -----------------------------------------------------------------------------
\subsection{Ideation and Brainstorming Process}

\begin{enumerate}[leftmargin=1.5em]
  \item \textbf{Problem Identification Phase:} Each team member individually
        used existing quiz platforms and documented pain points encountered as
        students.

  \item \textbf{Feature Brainstorming Sessions:} The team conducted multiple
        brainstorming sessions over Google Meet using collaborative tools:
        \begin{itemize}
          \item \textbf{Excalidraw} --- for quick whiteboard-style sketching of
                page layouts, user flows, and architecture diagrams.
          \item \textbf{Figma} --- for higher-fidelity wireframes and component
                design. Mockups of Landing Page, Dashboards, and Quiz Arena
                were created to align the team before coding.
        \end{itemize}

  \item \textbf{Research Phase:} Formal research on six competing platforms to
        identify market gaps, with a specific deep-dive on unauthorised access
        prevention.

  \item \textbf{Architecture Design:} Settled on a client-server architecture
        with clear separation between student and teacher interfaces.
\end{enumerate}

\textbf{Prioritised Feature Backlog:}

\begin{table}[H]
  \centering
  \caption{Feature Priority Classification}
  \small
  \begin{tabularx}{\linewidth}{@{} l X @{}}
    \toprule
    \textbf{Priority} & \textbf{Features} \\
    \midrule
    High & Student \& Teacher Dashboards, Create Quiz, Create Classroom, AI-generated questions, anti-cheat module, unauthorised-student detection \\
    Low  & Question tags, Mistake Book, detailed performance metrics, Coin/XP engagement mechanism \\
    \bottomrule
  \end{tabularx}
\end{table}

\begin{infobox}
  \textbf{Brainstorming Artifacts:} Excalidraw sketches, Figma wireframes, and
  Google Meet session screenshots are available in the project repository
  (\texttt{/assets/} folder) and can be provided on request.
\end{infobox}

% ── Brainstorming image placeholders ─────────────────────────────────────────

\begin{figure}[H]
  \centering
  \includegraphics[width=0.47\linewidth]{assets/excalidrawBrainStorm.png}
  \hfill
  \includegraphics[width=0.47\linewidth]{assets/figmaBrainstorm.png}
  \caption{Brainstorming Artifacts --- Excalidraw sketch (left) and Figma wireframe (right)}
\end{figure}

\begin{figure}[H]
  \centering
  \includegraphics[width=0.47\linewidth]{assets/googleMeet.jpeg}
  \caption{Team Brainstorming Session --- Google Meet}
\end{figure}

% -----------------------------------------------------------------------------
\subsection{Technology Stack and Justification}

\begin{table}[H]
  \centering
  \caption{Technology Stack with Justification}
  \small
  \begin{tabularx}{\linewidth}{@{} l l l X @{}}
    \toprule
    \textbf{Layer} & \textbf{Technology} & \textbf{Version} & \textbf{Why This Choice?} \\
    \midrule
    Frontend Framework & React        & 19.x & Component-based architecture ideal for a multi-page app; rich ecosystem; team familiarity \\
    Build Tool         & Vite         & 8.x  & Extremely fast dev server with HMR; near-instant rebuilds vs.\ Webpack \\
    Styling            & Tailwind CSS + Vanilla CSS & --- & Utility-first classes for rapid styling; Vanilla CSS for complex component-specific styles \\
    Backend Runtime    & Node.js      & LTS  & JavaScript across the full stack; non-blocking I/O suited for real-time quiz features \\
    Linting            & ESLint       & 9.x  & Early bug detection; enforces consistent code style across 6 team members \\
    \bottomrule
  \end{tabularx}
\end{table}

% -----------------------------------------------------------------------------
\subsection{Deployment Plan}

\begin{table}[H]
  \centering
  \caption{Live Deployment Plan}
  \small
  \begin{tabularx}{\linewidth}{@{} l l X @{}}
    \toprule
    \textbf{Stage} & \textbf{Tool / Service} & \textbf{Details} \\
    \midrule
    Version Control     & Git + GitHub    & Repository at \texttt{DevAggarwal03/WebTech\_Project} \\
    Frontend Hosting    & Vercel / Netlify & Free tier; automatic deployments from GitHub; global CDN \\
    Backend Hosting     & Render / Railway & Free tier; Node.js support; auto-deploy from GitHub \\
    Database            & MongoDB Atlas   & Free tier (512 MB); cloud-hosted; native JavaScript driver \\
    \bottomrule
  \end{tabularx}
\end{table}

% -----------------------------------------------------------------------------
\subsection{Data Flow Diagrams (DFD)}

\subsubsection{Level 0 --- Context Diagram}

\begin{figure}[H]
  \centering
  \includegraphics[width=0.85\linewidth]{assets/dfd0.png}
  \caption{Level 0 DFD — Context Diagram}
\end{figure}

\subsubsection{Level 1 --- Major Processes}

\begin{figure}[H]
  \centering
  \includegraphics[width=0.85\linewidth]{assets/dfd1.png}
  \caption{Level 1 DFD — Major Processes}
\end{figure}

% =============================================================================
\section{Implementation}
% =============================================================================

\subsection{Current Implementation Status ($\approx$50\%)}

The frontend of the application has been implemented with the following
completed pages and components:

\begin{table}[H]
  \centering
  \caption{Component Implementation Status}
  \small
  \begin{tabularx}{\linewidth}{@{} c l l c X @{}}
    \toprule
    \textbf{\#} & \textbf{Component} & \textbf{File} & \textbf{Status} & \textbf{Description} \\
    \midrule
    1 & Landing Page       & \texttt{QuizzLanding.jsx}     & \yes Complete & Hero section, feature cards, navbar, footer \\
    2 & Authentication     & \texttt{AuthPage.jsx}         & \yes Complete & Role toggle, sign-up form, animated background \\
    3 & Student Dashboard  & \texttt{StudentDashboard.jsx} & \yes Complete & Welcome banner, XP \& Rank, badges, quiz grid \\
    4 & Teacher Dashboard  & \texttt{TeacherDashboard.jsx} & \yes Complete & Stats banner, subject badges, quiz submissions \\
    5 & Quiz Arena         & \texttt{QuizArena.jsx}        & \yes Complete & 20 questions, timer, mark-for-review, calculator \\
    6 & Feature Card       & \texttt{FeatureCard.jsx}      & \yes Complete & Reusable landing page card component \\
    7 & Backend Server     & \texttt{backend/index.js}     & \no Scaffold  & Node.js initialised; implementation pending \\
    \bottomrule
  \end{tabularx}
\end{table}

% -----------------------------------------------------------------------------
\subsection{Key Implementation Highlights}

\subsubsection{Quiz Arena --- Full Interactive Quiz-Taking Experience}

The Quiz Arena is the most feature-rich component --- the dedicated area where
students take their quizzes. Key features include:

\begin{itemize}[leftmargin=1.5em]
  \item \textbf{Countdown Timer:} Real-time countdown from 9:35, with urgent
        styling when under 60 seconds remain.
  \item \textbf{Question Navigation:} Previous/Next buttons, question status
        grid for jumping to any question.
  \item \textbf{Mark for Review:} Students can flag questions to revisit, with
        a visual star indicator.
  \item \textbf{Colour-Coded Status Panel:} \textcolor{accent}{Green} =
        Answered, \textcolor{orange}{Yellow} = Marked for Review,
        \textcolor{gray}{Grey} = Not Visited.
  \item \textbf{Built-in Calculator:} Fully functional calculator for
        math-related quizzes, with basic arithmetic, clear, and backspace.
\end{itemize}

\subsubsection{Dual-Role Authentication}

The Auth page features:
\begin{itemize}[leftmargin=1.5em]
  \item Student/Teacher role toggle --- determines which dashboard the user
        sees after login.
  \item Animated binary-rain background for visual appeal.
  \item Form validation with name, email, and password fields.
\end{itemize}

\subsubsection{Dashboard Personalisation}

Both dashboards include:
\begin{itemize}[leftmargin=1.5em]
  \item Personalised welcome messages.
  \item Key stats at a glance (XP, Rank for students; Quiz count for teachers).
  \item Lightweight engagement: XP counter and achievement badges for students;
        subject badges for teachers.
  \item Recent quiz activity with relevant metadata.
\end{itemize}

\subsubsection{Component Architecture}

\begin{lstlisting}
frontend/src/
+-- components/
|   +-- AuthPage.jsx          (127 lines)
|   +-- FeatureCard.jsx       (reusable card)
|   +-- QuizArena.jsx         (221 lines -- largest component)
|   +-- QuizzLanding.jsx      (88 lines)
|   +-- StudentDashboard.jsx  (106 lines)
|   `-- TeacherDashboard.jsx  (118 lines)
+-- styles/
|   +-- AuthPage.css          (4.8 KB)
|   +-- FeatureCard.css       (0.8 KB)
|   +-- LandingPage.css       (0.7 KB)
|   +-- QuizArena.css         (7.6 KB -- most detailed)
|   +-- QuizzLanding.css      (3.1 KB)
|   +-- StudentDashboard.css  (7.2 KB)
|   `-- TeacherDashboard.css  (7.8 KB)
+-- App.jsx                   (main router)
+-- main.jsx                  (React entry point)
+-- index.css                 (global styles)
`-- App.css                   (app-level overrides)
\end{lstlisting}

% =============================================================================
\section{Future Work and Teamwork}
% =============================================================================

\subsection{Remaining Work and Next Steps}

\begin{longtable}{@{} c p{4.5cm} l p{4cm} @{}}
  \toprule
  \textbf{\#} & \textbf{Task} & \textbf{Priority} & \textbf{Scope} \\
  \midrule
  \endfirsthead
  \toprule
  \textbf{\#} & \textbf{Task} & \textbf{Priority} & \textbf{Scope} \\
  \midrule
  \endhead
  1 & Backend API --- Express.js REST API with JWT auth, quiz CRUD, classroom management & Critical & Backend \\
  2 & Database integration --- MongoDB schemas for Users, Quizzes, Classrooms, Scores & Critical & Backend \\
  3 & Unauthorised student detection --- Roster matching, approval queue, auto-expiring codes & High & Full Stack \\
  4 & Anti-cheat module --- Tab-switch detection, copy/paste blocking, timer enforcement & High & Frontend \\
  5 & AI-powered question generation --- Targeting each student's weak topics via AI APIs & High & Backend (API integration) \\
  6 & Student performance metrics --- Analytics, concept-wise tracking, strengths/weaknesses & Medium & Full Stack \\
  7 & Question tags --- Categorise by topic and difficulty for filtering and targeted practice & Medium & Full Stack \\
  8 & Mistake Book --- Personal error log for students to review wrong answers & Medium & Full Stack \\
  9 & Coin/XP mechanism --- Engagement system with XP earned from quiz participation & Low & Full Stack \\
  10 & Deployment --- Vercel/Netlify (frontend), Render (backend), MongoDB Atlas (DB) & Final & DevOps \\
  \bottomrule
  \caption{Task Backlog for Remaining Implementation}
\end{longtable}

% -----------------------------------------------------------------------------
\subsection{Integration Components}

\begin{table}[H]
  \centering
  \caption{Key Integration Components}
  \small
  \begin{tabularx}{\linewidth}{@{} l X @{}}
    \toprule
    \textbf{Component} & \textbf{Integration Details} \\
    \midrule
    Authentication         & JWT-based auth connecting frontend AuthPage to backend API \\
    Quiz Data Pipeline     & Quiz creation $\to$ DB storage $\to$ Quiz Arena rendering $\to$ Score submission \\
    Classroom Enrolment    & Join-code generation $\to$ Student request $\to$ Teacher approval $\to$ Roster sync \\
    AI Question Generation & Student performance data $\to$ AI API call $\to$ Personalised questions $\to$ Review centre \\
    Real-time Features     & WebSocket integration for live quiz sessions (potential) \\
    \bottomrule
  \end{tabularx}
\end{table}

% -----------------------------------------------------------------------------
\subsection{Team Contribution}

The team of 6 members is divided into two focused groups:

\begin{table}[H]
  \centering
  \caption{Team Roles and Responsibilities}
  \small
  \begin{tabularx}{\linewidth}{@{} l l X @{}}
    \toprule
    \textbf{Role} & \textbf{Members} & \textbf{Responsibilities} \\
    \midrule
    Research + Documentation & Dev, Mudit, Payal & Platform research, competitor analysis, feature specification, documentation, UI/UX design decisions, deployment planning \\
    Development              & Bhavya, Ayush, Rohan & Frontend component development, CSS styling, backend API development, database schema design, integration \\
    \bottomrule
  \end{tabularx}
\end{table}

\begin{infobox}
  All members participated in ideation discussions, feature prioritization, and
  code reviews.
\end{infobox}

% =============================================================================
\section{Additional Information}
% =============================================================================

\subsection{Key Differentiators from Existing Platforms}

\begin{table}[H]
  \centering
  \caption{Feature Comparison: QuizPod vs.\ Competitors}
  \small
  \begin{tabularx}{\linewidth}{@{} X c c c c @{}}
    \toprule
    \textbf{Feature} & \textbf{Kahoot} & \textbf{Quizlet} & \textbf{Google Forms} & \textbf{QuizPod} \\
    \midrule
    Teacher approval for enrolment  & \no & \no & \no & \yes \\
    Roster-based access control     & \no & \no & \no & \yes \\
    Auto-expiring join codes        & \no & \no & \no & \yes \\
    In-quiz calculator              & \no & \no & \no & \yes \\
    Role-based dashboards           & \no & \no & \no & \yes \\
    XP \& Badges                   & \yes (speed) & \no & \no & \yes (mastery) \\
    Free tier limitations           & 10 players & Paywalled & Basic only & \textbf{No paywall} \\
    Anti-cheat (tab detection)      & \no & \no & \no & \yes (planned) \\
    AI-generated questions          & \no & \no & \no & \yes (planned) \\
    Mistake Book                    & \no & \no & \no & \yes (planned) \\
    \bottomrule
  \end{tabularx}
\end{table}

% -----------------------------------------------------------------------------
\subsection{Project Repository}

\begin{itemize}[leftmargin=1.5em]
  \item \textbf{GitHub:} \href{https://github.com/DevAggarwal03/WebTech_Project}{\texttt{DevAggarwal03/WebTech\_Project}}
  \item \textbf{Structure:} Monorepo with \texttt{/frontend} (React + Vite) and \texttt{/backend} (Node.js)
\end{itemize}

% -----------------------------------------------------------------------------
\subsection{References}

\begin{enumerate}[leftmargin=1.5em]
  \item Kahoot --- \url{https://kahoot.com}
  \item Quizlet --- \url{https://quizlet.com}
  \item Google Forms --- \url{https://forms.google.com}
  \item Mentimeter --- \url{https://www.mentimeter.com}
  \item Socrative --- \url{https://www.socrative.com}
  \item Typeform --- \url{https://www.typeform.com}
  \item Google Classroom --- \url{https://classroom.google.com}
  \item Moodle LMS --- \url{https://moodle.org}
  \item Canvas LMS --- \url{https://www.instructure.com/canvas}
  \item Edmodo (discontinued 2022) --- \url{https://en.wikipedia.org/wiki/Edmodo}
  \item FlexiQuiz --- \url{https://www.flexiquiz.com}
  \item ClassMarker --- \url{https://www.classmarker.com}
  \item ProProfs Quiz Maker --- \url{https://www.proprofs.com/quiz-school}
\end{enumerate}

\vspace{1.5cm}
\noindent\rule{\linewidth}{0.4pt}\\[0.3em]

% ─────────────────────────────────────────────────────────────────────────────
\end{document}
% ─────────────────────────────────────────────────────────────────────────────
